% GENERATED FILE. Edit canonical lessons, then run npm run generate:books.
% canonical-source-hash: fnv1a64:419cfad38a8ed7bf
% canonical-lessons: RU-C444-ratusha, RU-C444-sobor, RU-C444-kolokolnya, RU-C444-krepost, RU-C444-meriya

\chapter{The old town}
\label{ch:ru-the-old-town}
\hlchaptermodality{444}

\begin{chapteropening}
\textbf{By the end of this chapter:} \emph{I can name the buildings of an old town: the town hall, the cathedral, the bell tower, the fortress and the city hall.}

Complete the last lesson of chapter 444: I can name the buildings of an old town: the town hall, the cathedral, the bell tower, the fortress and the city hall.
\end{chapteropening}

\section[\texorpdfstring{\ru{ратуша} (rátusha)}{ратуша (rátusha)}]{\ru{ратуша} (rátusha) — a town hall (in an old town)}

\label{lesson:RU-C444-ratusha}



\begin{quote}
Before the new one: say the Russian for a suburb, then the Russian for a roadway.
\end{quote}

\subsection*{You'll want to know: \ru{ратуша}}
\textbf{\ru{ратуша}} (\emph{rátusha}) — \textquotedblleft{}a town hall (in an old town)\textquotedblright{}.

\begin{quote}
\ru{На} \ru{площади} \ru{стоит} \ru{старая} \ru{ратуша}.
\end{quote}

\begin{quote}
\emph{Na plóshchadi stoít stáraya rátusha.}
\end{quote}

\begin{quote}
The old town hall stands on the square.
\end{quote}

Feminine. Said of old town halls in central Europe and the Baltic.

\begin{cousinweb}
\emph{\ru{ратуша}} came through Old East Slavic \emph{ratuša} from Polish \emph{ratusz} (EtymDB), a cousin of German \emph{Rathaus} 'a council house'.
\end{cousinweb}

\begin{grammarlens}[title={What you've built}]
One more word for telling, explaining and describing at B1.
\end{grammarlens}

\subsection*{Your turn}
\begin{itemize}
\raggedright
  \item \emph{Say it:} \emph{rátusha}
  \item \emph{Say it:} \emph{rátusha}, once more
  \item \emph{Say it:} a sentence of your own with \emph{rátusha}
  \item \emph{Recall:} say the Russian for a suburb, then the Russian for a roadway, then say \emph{rátusha} again
\end{itemize}

\begin{tcolorbox}[breakable,colback=teal!4,colframe=teal!35!black,title={Before you move on}]
What does \ru{ратуша} mean? (\textquotedblleft{}A town hall (in an old town)\textquotedblright{}.) Say it once more.
\end{tcolorbox}
\section[\texorpdfstring{\ru{собор} (sobór)}{собор (sobór)}]{\ru{собор} (sobór) — a cathedral; a council}

\label{lesson:RU-C444-sobor}



\begin{quote}
Before the new one: say the Russian for a roadway, then the Russian for a town hall.
\end{quote}

\subsection*{You'll want to know: \ru{собор}}
\textbf{\ru{собор}} (\emph{sobór}) — \textquotedblleft{}a cathedral; a council\textquotedblright{}.

\begin{quote}
\ru{Рядом} \ru{с} \ru{ратушей} — \ru{большой} \ru{собор}.
\end{quote}

\begin{quote}
\emph{Ryádom s rátushey — bolshóy sobór.}
\end{quote}

\begin{quote}
Next to the town hall there is a large cathedral.
\end{quote}

Masculine. A big main church; also an assembly.

\begin{cousinweb}
\emph{\ru{собор}} came from Church Slavonic \emph{sŭborŭ} 'a meeting, a council' (EtymDB), which is inherited from Proto-Slavic \emph{sŭborŭ} 'a gathering', a noun usually linked with \emph{\ru{собирать}} 'to gather': first the gathering of the faithful, then their great church.
\end{cousinweb}

\begin{grammarlens}[title={What you've built}]
One more word for telling, explaining and describing at B1.
\end{grammarlens}

\subsection*{Your turn}
\begin{itemize}
\raggedright
  \item \emph{Say it:} \emph{sobór}
  \item \emph{Say it:} \emph{sobór}, once more
  \item \emph{Say it:} a sentence of your own with \emph{sobór}
  \item \emph{Recall:} say the Russian for a roadway, then the Russian for a town hall, then say \emph{sobór} again
\end{itemize}

\begin{tcolorbox}[breakable,colback=teal!4,colframe=teal!35!black,title={Before you move on}]
What does \ru{собор} mean? (\textquotedblleft{}A cathedral; a council\textquotedblright{}.) Say it once more.
\end{tcolorbox}
\section[\texorpdfstring{\ru{колокольня} (kolokólnya)}{колокольня (kolokólnya)}]{\ru{колокольня} (kolokólnya) — a bell tower}

\label{lesson:RU-C444-kolokolnya}



\begin{quote}
Before the new one: say the Russian for a town hall, then the Russian for a cathedral.
\end{quote}

\subsection*{You'll want to know: \ru{колокольня}}
\textbf{\ru{колокольня}} (\emph{kolokólnya}) — \textquotedblleft{}a bell tower\textquotedblright{}.

\begin{quote}
\ru{С} \ru{колокольни} \ru{виден} \ru{весь} \ru{город}.
\end{quote}

\begin{quote}
\emph{S kolokólni víden ves górod.}
\end{quote}

\begin{quote}
The whole town can be seen from the bell tower.
\end{quote}

Feminine. Built on \emph{\ru{колокол}} — a bell.

\begin{cousinweb}
\emph{\ru{колокольня}} is built from \emph{\ru{колокол}} 'a bell' and \emph{-\ru{ня}}, an ending for places (EtymDB). \emph{\ru{Колокол}} is inherited from Proto-Slavic \emph{kolkolŭ} 'a bell' (EtymDB).
\end{cousinweb}

\begin{grammarlens}[title={What you've built}]
One more word for telling, explaining and describing at B1.
\end{grammarlens}

\subsection*{Your turn}
\begin{itemize}
\raggedright
  \item \emph{Say it:} \emph{kolokólnya}
  \item \emph{Say it:} \emph{kolokólnya}, once more
  \item \emph{Say it:} a sentence of your own with \emph{kolokólnya}
  \item \emph{Recall:} say the Russian for a town hall, then the Russian for a cathedral, then say \emph{kolokólnya} again
\end{itemize}

\begin{tcolorbox}[breakable,colback=teal!4,colframe=teal!35!black,title={Before you move on}]
What does \ru{колокольня} mean? (\textquotedblleft{}A bell tower\textquotedblright{}.) Say it once more.
\end{tcolorbox}
\section[\texorpdfstring{\ru{крепость} (krépost)}{крепость (krépost)}]{\ru{крепость} (krépost) — a fortress}

\label{lesson:RU-C444-krepost}



\begin{quote}
Before the new one: say the Russian for a cathedral, then the Russian for a bell tower.
\end{quote}

\subsection*{You'll want to know: \ru{крепость}}
\textbf{\ru{крепость}} (\emph{krépost}) — \textquotedblleft{}a fortress\textquotedblright{}.

\begin{quote}
\ru{На} \ru{холме} \ru{стоит} \ru{старая} \ru{крепость}.
\end{quote}

\begin{quote}
\emph{Na kholmé stoít stáraya krépost.}
\end{quote}

\begin{quote}
An old fortress stands on the hill.
\end{quote}

Feminine. Also strength: \emph{\ru{крепость} \ru{чая}} — the strength of tea.

\begin{cousinweb}
\emph{\ru{крепость}} is inherited through Old East Slavic \emph{krěpostĭ} 'strength' from Proto-Slavic \emph{krěpostĭ} (EtymDB), on the root of \emph{\ru{крепкий}} 'strong': first strength, then a strong place.
\end{cousinweb}

\begin{grammarlens}[title={What you've built}]
One more word for telling, explaining and describing at B1.
\end{grammarlens}

\subsection*{Your turn}
\begin{itemize}
\raggedright
  \item \emph{Say it:} \emph{krépost}
  \item \emph{Say it:} \emph{krépost}, once more
  \item \emph{Say it:} a sentence of your own with \emph{krépost}
  \item \emph{Recall:} say the Russian for a cathedral, then the Russian for a bell tower, then say \emph{krépost} again
\end{itemize}

\begin{tcolorbox}[breakable,colback=teal!4,colframe=teal!35!black,title={Before you move on}]
What does \ru{крепость} mean? (\textquotedblleft{}A fortress\textquotedblright{}.) Say it once more.
\end{tcolorbox}
\section[\texorpdfstring{\ru{мэрия} (mériya)}{мэрия (mériya)}]{\ru{мэрия} (mériya) — a city hall, a mayor's office}

\label{lesson:RU-C444-meriya}



\begin{quote}
Before the new one: say the Russian for a bell tower, then the Russian for a fortress.
\end{quote}

\subsection*{You'll want to know: \ru{мэрия}}
\textbf{\ru{мэрия}} (\emph{mériya}) — \textquotedblleft{}a city hall, a mayor's office\textquotedblright{}.

\begin{quote}
\ru{Документы} \ru{можно} \ru{получить} \ru{в} \ru{мэрии}.
\end{quote}

\begin{quote}
\emph{Dokuménty mózhno poluchít v mérii.}
\end{quote}

\begin{quote}
You can get the documents at the city hall.
\end{quote}

Feminine. \emph{\ru{Мэр}} — a mayor.

\begin{cousinweb}
\emph{\ru{мэрия}} was borrowed from French \emph{mairie} 'a mayor's office' (EtymDB), built on \emph{maire} 'a mayor'.
\end{cousinweb}

\begin{grammarlens}[title={What you've built}]
One more word for telling, explaining and describing at B1.
\end{grammarlens}

\subsection*{Your turn}
\begin{itemize}
\raggedright
  \item \emph{Say it:} \emph{mériya}
  \item \emph{Say it:} \emph{mériya}, once more
  \item \emph{Say it:} a sentence of your own with \emph{mériya}
  \item \emph{Recall:} say the Russian for a bell tower, then the Russian for a fortress, then say \emph{mériya} again
\end{itemize}

\begin{tcolorbox}[breakable,colback=teal!4,colframe=teal!35!black,title={Before you move on}]
What does \ru{мэрия} mean? (\textquotedblleft{}A city hall, a mayor's office\textquotedblright{}.) Say it once more.
\end{tcolorbox}
